\documentclass{article}

\usepackage{Paper}
\pdftitle{Block-E4-PDP2}
\usepackage{listings}
\lstset{
    basicstyle=\ttfamily\normalsize,
    keywordstyle=\color{blue}\bfseries,
    commentstyle=\color{green!60!black},
    stringstyle=\color{red},
    showstringspaces=false,
    numbers=left,
    numberstyle=\tiny,
    frame=single,
    breaklines=true
}

\author{Matteo Frongillo}

\begin{document}
\pagestyle{plain}

\subsection*{Scaffold}
\begin{lstlisting}[language=C++]
// ===================== INCLUDES =====================
#include "Seeed_AMG8833_driver.h"
#include <Grove_I2C_Motor_Driver.h>
#include <Wire.h>
#include <stdint.h>
#include <math.h>

// ===================== CONSTANTS / GLOBALS =====================

// Motor driver settings
#define I2C_ADDRESS 0x0f

// --- CALIBRATION VALUES ---
// Replace these only after validation
float ambientTempC = AMBIENT_TEMP_C;
float boardSideCm = BOARD_SIDE_CM;
char prototypePlacement[] = PLACEMENT; // CORNER or BORDER

// --- MOTOR CONFIGURATION ---
// Keep this style similar to the original sketch
const int START_OFFSET_STEPS = 0;
int SWEEP_STEPS = 128; // update according to placement
int currentMotorPosition = 0;

// Approximate conversion between steps and angle
// AI may adjust this if needed but should keep the same idea
float stepsPerDegree = 128.0 / 90.0;

// Ultrasonic sensor pin (Grove)
const int ultrasonicPin = 5;

// Thermal camera
AMG8833 sensor;

#define ROWS 8
#define COLS 8
#define TOTAL_PIXELS 64

// ===================== HELPER FUNCTIONS =====================

// --- REQUIRED THRESHOLD FUNCTION ---
float getHotThreshold(float distanceCm, float boardSideLength, float ambientTempC) {
  if (distanceCm <= 0.45 * boardSideLength) {
    return 30.0;
  } else if (distanceCm < 0.9 * boardSideLength) {
    return ambientTempC + 38.72 * exp(-0.04661 * distanceCm);
  } else {
    return max(24.5, ambientTempC + 1.0);
  }
}

// --- MOTOR / ANGLE HELPERS ---
// TODO:
// - set SWEEP_STEPS from placement
// - convert motor position to angle
// - keep 0° -> X° -> 0° sweep behavior

// --- ULTRASONIC HELPERS ---
long readUltrasonicDuration() {
  pinMode(ultrasonicPin, OUTPUT);
  digitalWrite(ultrasonicPin, LOW);
  delayMicroseconds(2);
  digitalWrite(ultrasonicPin, HIGH);
  delayMicroseconds(10);
  digitalWrite(ultrasonicPin, LOW);

  pinMode(ultrasonicPin, INPUT);
  return pulseIn(ultrasonicPin, HIGH, 30000);
}

float microsecondsToCentimeters(long duration) {
  if (duration == 0) return -1.0;
  return (duration * 0.0343) / 2.0;
}

// TODO:
// - optionally average multiple readings
// - return distance in cm
// - apply valid-range rule 5 < x < d in detection logic

// --- THERMAL HELPERS ---
// TODO:
// - read all 64 pixels
// - compute max temperature over the full matrix
// - do not use only middle columns

// --- DETECTION / REPORTING HELPERS ---
// TODO:
// - combine valid distance + thermal threshold condition
// - print:
//   current distance
//   current motor angle
//   maximum temperature
//   detection status

// ===================== SETUP =====================
void setup() {
  Serial.begin(115200);

  Wire.begin();
  Motor.begin(I2C_ADDRESS);

  Motor.stop(MOTOR1);
  Motor.stop(MOTOR2);

  sensor.init();

  delay(1000);

  // TODO:
  // - choose SWEEP_STEPS based on placement
  // - optionally move to starting offset
  if (START_OFFSET_STEPS != 0) {
    int direction = (START_OFFSET_STEPS > 0) ? 1 : -1;
    int stepsToMove = abs(START_OFFSET_STEPS);

    for (int i = 0; i < stepsToMove; i++) {
      Motor.StepperRun(direction, 0, 0);
      delay(10);
    }
    currentMotorPosition = START_OFFSET_STEPS;
  }
}

// ===================== MAIN LOOP =====================
void loop() {
  currentMotorPosition = 0;

  delay(500);

  // ---------- Sweep Forward 0 -> X ----------
  for (int i = 0; i < SWEEP_STEPS; i++) {
    Motor.StepperRun(1, 0, 0);
    currentMotorPosition++;

    // TODO:
    // - compute current angle in degrees
    // - read ultrasonic sensor
    // - read all thermal pixels and max temperature
    // - compute threshold
    // - detect only if:
    //   valid distance AND maxTemp > threshold
    // - print serial status

    delayMicroseconds(5000);
  }

  delay(2000);

  // ---------- Sweep Backward X -> 0 ----------
  for (int i = 0; i < SWEEP_STEPS; i++) {
    Motor.StepperRun(-1, 0, 0);
    currentMotorPosition--;

    // TODO:
    // - same logic as forward sweep
    // - keep angle tracking consistent

    delayMicroseconds(5000);
  }

  delay(2000);
}
\end{lstlisting}

\newpage
\section*{Prompt blocks}
\subsection*{Block A: Includes}
\begin{lstlisting}
You are editing an Arduino Uno `.ino` sketch step by step.

RULES
- Keep the code very similar to the provided sketch in structure, naming style, and general style.
- Do not redesign the project architecture.
- Keep the section order:
  1. includes
  2. constants / globals
  3. helper functions
  4. setup
  5. loop
- Keep the sketch compilable after this step.
- Return the full updated sketch.
- Modify only what is needed for this step.

TASK FOR THIS STEP
Review only the INCLUDES section.

Use these includes exactly:
- #include "Seeed_AMG8833_driver.h"
- #include <Grove_I2C_Motor_Driver.h>
- #include <Wire.h>
- #include <stdint.h>

You may add only strictly necessary standard Arduino/math includes if required for compilation.
Do not change the rest of the sketch unless strictly necessary.

SECTION TO REVIEW
// ===================== INCLUDES =====================
#include "Seeed_AMG8833_driver.h"
#include <Grove_I2C_Motor_Driver.h>
#include <Wire.h>
#include <stdint.h>
#include <math.h>

CURRENT FULL CODE
[PASTE THE INITIAL SCAFFOLD HERE]
\end{lstlisting}


\end{document}